% Source: https://artofproblemsolving.com/community/c3h2213845_beginner_problems
% Stable identity: aops-topic-2213845-q03
\textbf{The configuration.} A side of the rectangle containing two vertices of the hexagon contains the segment joining them, and since the hexagon lies inside the rectangle, this segment must be a side of the hexagon. With six vertices on four sides, at least two sides of the rectangle contain two vertices each, hence contain hexagon sides. Two sides of a regular hexagon are never perpendicular (their directions differ by multiples of $60^\circ$), so these rectangle sides are opposite, and they contain two opposite sides of the hexagon. The other two vertices lie on the other two sides of the rectangle.

\begin{asy}
usepackage("amsmath");
size(5cm);
defaultpen(fontsize(10pt));
real s=2; pair[] V; for(int k=0;k<6;++k) V.push(s*dir(60*k));
draw((-s,-s*sqrt(3)/2)--(s,-s*sqrt(3)/2)--(s,s*sqrt(3)/2)--(-s,s*sqrt(3)/2)--cycle,blue+linewidth(0.9));
draw(V[0]--V[1]--V[2]--V[3]--V[4]--V[5]--cycle,linewidth(0.9));
draw(V[0]--V[3],gray+dashed); draw(V[1]--V[5],red);
for(int k=0;k<6;++k) dot(V[k]);
label("$2s$",(-0.6,0),N,gray); label("$s\sqrt3$",(V[1]+V[5])/2+(0,0.55),W,red);
label("$s$",(V[1]+V[2])/2,N);
\end{asy}

\textbf{The area.} If the hexagon has side $s$, the rectangle has width $2s$ (the long diagonal) and height $s\sqrt3$ (the distance between opposite sides), so its area is $2\sqrt3\,s^2$. The hexagon consists of six equilateral triangles of side $s$, so its area is $\dfrac{3\sqrt3}{2}s^2$. Hence
\[\dfrac{[\text{rectangle}]}{[\text{hexagon}]}=\dfrac{2\sqrt3\,s^2}{\dfrac{3\sqrt3}{2}s^2}=\dfrac{4}{3},\qquad[\text{rectangle}]=\dfrac{4}{3}\cdot24=\boxed{32}.\]
